\documentclass[conference]{IEEEtran}
\usepackage{graphicx}
\usepackage{booktabs}
\usepackage{url}
\usepackage{amsmath}
\title{PIVOT: An Auditable CHIA Loop for Near-Memory Binary-Vector Design}
\author{\IEEEauthorblockN{Author names to be supplied by the team}
\IEEEauthorblockA{Affiliation and contact to be supplied}}
\begin{document}
\maketitle
\begin{abstract}
PIVOT explores the design of a near-memory binary-vector tile using a frozen
functional evaluator, restricted address-mapping programs, and matched
grid/random controls. Hyperdimensional operations and a simplified spatial
pooler motivate shared bitwise, population-count, selection, and saturation
primitives. The artifact includes a Python reference, parameterized RTL and
Chisel, a differential verification harness, a cycle-approximate memory model,
and composable CHIA evaluation nodes. An optional API-only language-model
proposer receives aggregate feedback without access to the reference or test
generator. The current local evidence is a reproducible synthetic model study;
distributed agent execution, simulator correctness, and physical PPA require
separate recorded runs. We present the evaluator and control protocol without
claiming unmeasured hardware improvements or agent superiority.
\end{abstract}

\section{Motivation and scope}
The target is logic near memory, rather than an HBM controller or physical
interface. HDC workloads motivate XOR binding, thresholded bundling, cyclic
permutation, and Hamming similarity. Sparse binary-vector workloads also need
overlap, selection, and saturating permanence updates. Sharing primitive
operations suggests a compact tile, but it does not establish that an agent
is needed to discover its implementation.

Our research question is whether bounded program synthesis of address and
issue policies improves on finite parameter search under a shared evaluator.
The controls are essential: a favorable comparison against one starting design
does not establish an advantage over random or exhaustive search. The artifact
therefore preserves both successful and rejected evaluations, records provenance,
and labels model estimates separately from physical measurements.

We deliberately exclude model fine-tuning, physical HBM implementation, FPGA
bring-up, and place-and-route. The current implementation also excludes
automatic top-k RTL rewriting and full Chipyard integration. Those boundaries
are visible in machine-readable status fields rather than hidden behind a
single successful workflow exit code.

\section{Tile and functional contract}
The legacy tile accepts two binary vectors, a packed candidate-vector set,
and an opcode. Candidate zero occupies the least significant packed bits.
The retained opcodes implement XOR, rotate-left, threshold vote, per-bit count,
Hamming distance, and nearest-K indices. Equal distances select the lower
candidate index. Unused output buses and reserved opcodes return zero. The
baseline is combinational; clock/reset ports do not imply a pipeline.

Parameter repairs cover one-element candidate sets, non-power-of-two rotate
widths, auxiliary-bus capacity, and invalid K values. A separate saturation
block implements permanence updates; the vote-counter opcode must not be
misrepresented as learning. The Chisel source mirrors these contracts.

The Python reference additionally provides HDC encode/query composition and
global spatial-pooler inhibition over supplied connectivity masks. These
algorithms are not a claim of a trained classifier or a complete HTM model.
Synthetic examples include reference answers and explicit provenance labels.

\section{Evaluator and isolation}
Directed cases cover zero/all-one patterns, thresholds, rotations, duplicate
candidates, stable ties, and boundary counts. Seeded random cases use separate
development and heldout streams. The HDL process receives inputs only; expected
answers remain in the Python comparator. The comparator checks order, every
output bus, unknown values, missing observations, and the completion marker.
Reference-source hashes detect accidental evaluator edits. Hash checks alone
are not an operating-system sandbox for malicious HDL.

The actual automatic proposal interface accepts JSON and bounded integer
expressions. It provides no shell, file reader, code-execution tool, or evaluator
source to the model. An AST interpreter rejects calls, imports, attributes,
unbounded shifts, and excessive expression depth. Every address mapping must be
injective over the complete benchmark burst domain and remain within physical
coordinate bounds. The evaluator fixes the workload and request count; candidates
cannot remove traffic. This restriction gives an enforceable interface boundary
for the implemented search domain, rather than claiming arbitrary RTL isolation.

\section{Memory model and CHIA cascade}
The read-only memory model exposes pseudo channels, banks, rows, burst transfer,
row hits, and row conflicts. Timing parameters reference gem5's HBM2 interface
\cite{gem5hbm}. This is parameterization, not calibration: refresh, bank-group
timing, command-bus details, writes, and physical effects remain omitted. Memory
and logical datapath cycles are added conservatively without overlap. Every
result is marked \emph{cycle-approximate} and \emph{uncalibrated} until compared
against external gem5 observations.

CHIA \cite{chia} supplies task dependencies and heterogeneous worker resource
tags. Tier 0 smoke-checks the shared RTL baseline; Tier 1 evaluates the candidate
and records gate/model outcomes; Tier 2 optionally invokes configured synthesis
tools. Artifacts return to the head before an ephemeral worker disappears. The
current synthesis node measures the shared baseline, not a distinct hardware
implementation for every logical lane-width candidate. Tier 3 is an explicit
unimplemented extension point. The local fallback calls the same core evaluator
without claiming that a distributed CHIA run occurred.

Tailscale configurations distinguish existing routed interfaces from CHIA's
userspace SOCKS/CONNECT relay. A managed GCP template is optional and never
launched by the default script. Result archiving precedes teardown because
CHIA-managed cloud instances can be deleted by \texttt{chia down}.

\section{Experimental protocol and current evidence}
The finite control space contains 144 combinations: four logical lane widths,
six prefetch depths, three layouts, and two issue policies. Grid search
enumerates this space; random search permutes it without replacement. Both
receive the same workload and attempted-evaluation budget. Invalid agent
proposals consume proposal budget. Wall time is recorded separately, so equal
evaluation count is not called equal monetary cost.

\begin{table}[t]
\centering
\caption{Evidence categories; update only from archived run manifests.}
\begin{tabular}{ll}
\toprule
Category & Initial artifact state\\
\midrule
Python reference / invariant tests & Executed locally\\
Synthetic examples & Generated, labeled synthetic\\
Grid/random memory estimates & Executed locally\\
Live Gemini proposal comparison & Requires private run\\
RTL differential 100k gate & Requires simulator report\\
Candidate physical area/frequency & Not measured\\
gem5 calibration & Not performed\\
\bottomrule
\end{tabular}
\end{table}

We executed three 100-attempt-per-arm model comparisons and one exhaustive
144-attempt-per-arm comparison, each with a shared baseline outside the arm
budget. The archive contains 888 search attempts and four baselines: 892
analytical evaluations, not 892 distinct chips or implemented candidates.
All records have status \texttt{model\_evaluated}; none has an RTL pass, an LLM
proposal, or a physical synthesis result. The evaluated source digest begins
\texttt{6107b2d7f979dbbb}; the complete digest, workload hashes, and checked CSV
hashes are recorded in \texttt{docs/LOCAL\_RESULTS.md}.

\begin{table}[t]
\centering
\caption{Actual local analytical results: best modeled operations per cycle.
The last row exhausts the complete finite control space.}
\begin{tabular}{rrrrr}
\toprule
Seed & Budget/arm & Baseline & Grid & Random\\
\midrule
20260923 & 100 & 0.095352 & 0.171951 & 0.316518\\
20260924 & 100 & 0.095298 & 0.171582 & 0.316440\\
20260925 & 100 & 0.095302 & 0.171731 & 0.316127\\
20260923 & 144 & 0.095352 & 0.316518 & 0.316518\\
\bottomrule
\end{tabular}
\end{table}

The first 100 grid configurations are a prefix of a fixed ordering that reaches
larger prefetch depths late. Random permutations examine a different subset;
their higher best estimates in these three truncated runs therefore do not
establish general search superiority. Both exhaustive arms reach exactly the
same optimum. No agent comparison is possible without recorded provider calls.
All candidates retain 2,560 logical operations and 5,120 memory bursts per trace.
The results are uncalibrated model estimates, not measured hardware speedups.

The artifact also contains 256 synthetic HDC and 256 synthetic HTM examples
with reference outputs and verified content hashes. No fitting or fine-tuning
was performed. Local validation comprises 39 core Python tests, 12 workflow
tests, and two Chisel-adapter tests. Chisel compilation and four CHIRRTL elaborations succeeded, including
default, one-bit, irregular-width, and standalone-permanence configurations.
These software checks and elaborations do not substitute for the pending RTL
differential runs. No historical candidate count or spend from the planning
document is treated as an experimental observation.

The initial model Pareto set minimizes modeled cycles, queue storage, and lane
width. It is not a silicon PPA front. The intended throughput-per-area metric
requires a verified candidate-specific implementation, a measured synthesis
frequency, and a reported area under stated constraints. Substituting generic
cell counts or shared-baseline area would conflate incompatible evidence.

\section{Limitations and next measurements}
The first required follow-up is functional simulation of both emitted Chisel
and legacy RTL, including permanence. Next, a real CHIA/Gemini run can compare
the restricted program domain against both controls across repeated seeds.
Physical claims require implementing selected policies in hardware, reverifying
them, and using consistent synthesis corners and timing constraints. External
gem5 traces are needed to assess the memory model's approximation error.

These omissions limit the present contribution to a reproducible evaluator,
control protocol, and reusable workflow artifact. They do not invalidate those
components, but they prevent a claim of a completed physical-design study.
Positive and negative agent results are both acceptable outcomes under the
same predeclared evaluator; neither can be substituted by unexecuted stages.

\section*{Reproducibility and AI assistance}
The artifact includes launch commands, source hashes, SQLite records, synthetic
data provenance, and explicit missing-tool statuses. Team members must review
this manuscript, supply author information, verify the final PDF length, and
publish the repository/result release. OpenAI Codex assisted implementation,
tests, documentation, and this manuscript draft. Human authors are responsible
for the submitted claims and evidence.

\begin{thebibliography}{9}
\bibitem{chia} CHIA authors, ``CHIA: An Open-Source Framework for Principled,
Agentic AI-Driven Hardware/Software Co-Design Research,'' arXiv:2606.27350, 2026.
\url{https://arxiv.org/abs/2606.27350}.
\bibitem{gem5hbm} gem5 project, ``HBM memory interfaces,''
\url{https://github.com/gem5/gem5/blob/stable/src/python/gem5/components/memory/dram_interfaces/hbm.py}.
\bibitem{hackathon} A$^3$ workshop, ``CHIA Hackathon,''
\url{https://agentic-arch.org/hackathon.html}.
\end{thebibliography}
\end{document}
